\documentclass[10pt,a4paper]{amsart}
\usepackage[margin=19mm]{geometry}
\usepackage{amsmath,amssymb,mathtools}
\usepackage[scr=rsfs,cal=euler]{mathalfa}
\usepackage{xfrac}
\usepackage[unicode,hidelinks,bookmarks=false]{hyperref}
\newcommand{\ie}{i.e.\ }
\newcommand{\cf}{cf.\ }
\newcommand{\resp}{respectively\ }
\newcommand{\Subsection}{\subsection}
\providecommand{\nobreakdash}{\nobreak\hbox{-}}
\setlength{\emergencystretch}{2em}
\multlinegap0pt
\usepackage{bm}
\usepackage{bbm}
\usepackage{dsfont}
\usepackage{xspace}
\usepackage{cancel}
\usepackage{mathtools}
\usepackage{amsthm}
\makeatletter
\@ifundefined{thm}{\newtheorem{thm}{Thm}}{}
\@ifundefined{prop}{\newtheorem{prop}[thm]{Proposition}}{}
\makeatother
\makeatletter
\newcommand{\cS}{{\mathcal S}}
\newcommand{\disc}{\operatorname{Disc}}
\expandafter\ifx\csname pfof\endcsname\relax
\newenvironment{pfof}[1]{\vspace{1ex}\noindent{\bf Proof of
#1}\hspace{0.5em}}{\hfill\qed\vspace{1ex}}
\fi
\newcommand{\tD}{{\widetilde D}}
\newcommand{\tU}{{\widetilde U}}
\newcommand{\tcS}{{\widetilde{\mathcal S}}}
\newcommand{\te}{{\tilde e}}
\makeatother
\title[Convergence to decorated L\'evy processes in non-Skorohod to]{Convergence to decorated L\'evy processes in non-Skorohod topologies for dynamical systems}
\author{Ana Cristina Moreira Freitas, Jorge Milhazes Freitas, Ian Melbourne, Mike Todd}
\date{Source: arXiv:2310.00978v3 (CC BY 4.0)}
\begin{document}
\maketitle

\noindent\emph{Source and licence.} Convergence to decorated L\'evy processes in non-Skorohod topologies for dynamical systems. Version arXiv:2310.00978v3 (CC BY 4.0), distributed under the Creative Commons Attribution 4.0 International licence, as recorded in the arXiv licence metadata for this exact version and shown on the abstract page https://arxiv.org/abs/2310.00978v3. Full rights notice: licenses/arxiv-2310.00978-CC-BY-4.0.txt.\par\smallskip

\noindent\textit{Editorial scope.} Counted unit: the Prop whose source label is \texttt{prop:comp2} in the article (source0.tex lines 1284--1286, source label \texttt{prop:comp2}), delivered together with the article's own complete original proof of it (source0.tex lines 1288--1377, one \texttt{proof} environment of 90 lines). No mathematics is written, repaired or re-derived: statement and proof are the article's own text line for line. Required context retained uncounted: eq:main (lines 896--898). Every definition, display and label of those lines is retained unchanged. Omitted from this package and not redistributed: 1--895, 899--1283, 1287--1287, 1378--2048. Nothing inside a retained excerpt is omitted or altered: every retained body is delivered exactly as the article's lines are written, so no omission receipt of any kind accompanies this package. Citations: the retained excerpts carry no citation command at all, so no bibliography entry of the article is transcribed and no reference is redistributed. Reference adaptation: none was needed and none was made; every reference inside the delivered unit resolves inside it, so no unresolved reference or citation is produced. Printed slips kept literal: no printed word, symbol, hypothesis or number of the counted statement or of its proof was altered. Layout and class: the article's own class is \texttt{article} with packages including graphicx, amsmath, amssymb, theorem, dutchcal, tikz, hyperref; the wrapper is the lane's pinned \texttt{amsart} editorial wrapper at 10pt on A4, which restates the theorem-like environments the retained text needs and adds the article's own macro definitions that the retained text needs (\texttt{\textbackslash cS}, \texttt{\textbackslash disc}, \texttt{\textbackslash tD}, \texttt{\textbackslash tU}, \texttt{\textbackslash tcS}, \texttt{\textbackslash te}), with extra wrapper packages (bm, bbm, dsfont, xspace, cancel, mathtools). The delivered numbering is the unit's own. Assets: no retained span contains a figure environment, a graphics-inclusion command, a PDF-page inclusion, a table or a TikZ picture; no image file is copied into this package, the package declares no asset dependency and no image file is redistributed with it. Licence: version arXiv:2310.00978v3 of this article is distributed under the Creative Commons Attribution 4.0 International licence, as recorded in the arXiv licence metadata for this exact version and shown on the abstract page https://arxiv.org/abs/2310.00978v3. A full rights notice is delivered at licenses/arxiv-2310.00978-CC-BY-4.0.txt. Characters: every retained excerpt is pure ASCII, so the delivered engine, XeLaTeX with its default Latin Modern text font, reports no missing character.
	\begin{equation} \label{eq:main}
	b_n^{-1}\max_{0\le j\le n} d_\tD(\xi,\zeta)\circ f^j \to_{\mu_X} 0.
	\end{equation} 

\begin{prop}  \label{prop:comp2}
$d_{F',[0,t_{n,N_n}]}(\chi U_n, \tU_n)\to_{\mu_X}0$.
\end{prop}

\begin{proof}
We claim that
\[
d_{F',[0,t_{n,N_n}]}(\chi U_n,\tU_n)
\le b_n^{-1}\max_{0\le j\le n}d_\tD(\xi,\zeta)\circ f^j.
\]
The result then follows by hypothesis~\eqref{eq:main}.

It remains to prove the claim.
In general, $\disc_{U_n}\subset\tcS_n=\{t_{n,j}:1\le j\le N_n\}$.
We first consider the slightly simpler case 
$\disc_{U_n}=\tcS_n$ for all $n$. Then
\[
\chi U_n=(U_n,\tcS_n,\{e_{U_n}^\tau\}_{\tau\in \cS_{U_n}}), \qquad
\tU_n=(U_n,\tcS_n,\{\te_{U_n}^\tau\}_{\tau\in \cS_{U_n}}),
\]
where 
\begin{align*}
e_{U_n}^\tau(t) & = U_n(t_{n,j})+ |\Delta U_n(\tau)| \Pi(\Delta U_n(\tau))(t)  + t\big\{\Delta U_n(\tau) -|\Delta U_n(\tau)|\Pi(\Delta U_n(\tau))(1) \big\},
\\
\te_{U_n}^\tau(t) & =U_n(t_{n,j})+b_n^{-1}v_{[tR]}\circ f^j,
\end{align*}
for $\tau=t_{n,j+1}$.

Note that $\Delta U_n(\tau)= b_n^{-1}V\circ f^j$.
Hence
\begin{align*}
e_{U_n}^\tau(t) & 
=U_n(t_{n,j})+b_n^{-1}\big\{
|V| \Pi(V)(t)  + t\big\{V-|V|\Pi(V)(1)
\big\} \circ f^j
\\ & =U_n(t_{n,j})+b_n^{-1} \zeta(t)\circ f^j.
\end{align*}
Also,
\[
\te_{U_n}^\tau(t)  =U_n(t_{n,j})+b_n^{-1}\xi(t)\circ f^j.
\]

On the interval $(t_{n,j},t_{n,j+1}]$, it follows that
$\pi_\tD \tU_n$ is the concatenation of the constant function $U_n(t_{n,j})$ with $U_n(t_{n,j})+b_n^{-1} \xi(t) \circ f^j$, while
$\pi_\tD \chi U_n$ is the concatenation of the constant function $U_n(t_{n,j})$ with $U_n(t_{n,j})+b_n^{-1}\zeta(t)\circ f^j$. 
Hence 
\[
d_{\tD,(t_{n,j},t_{n,j+1}]}(\pi_\tD \tU_n,\pi_\tD \chi U_n)
\le b_n^{-1}d_\tD(\xi,\zeta)\circ f^j
\]
and it follows that
\[
d_{\tD,[0,t_{n,N_n}]}(\pi_\tD \tU_n,\pi_\tD \chi U_n)
\le b_n^{-1}\max_{0\le j\le n}d_\tD(\xi,\zeta)\circ f^j.
\]

Also,
$d_{E,(t_{n,j},t_{n,j+1}]}(\pi_E \tU_n,\pi_E \chi U_n)=b_n^{-1}H\circ f^j$
where $H$ is the Hausdorff distance between the smallest closed boxes containing 
the ranges of the functions $\xi(t)$
and $\zeta(t)$ for $t\in[0,1]$.
(So $H=\max_{1\le k\le d}H_k$ where $H_k$ is the Hausdorff distance between 
the smallest closed interval containing 
$\{p_k\xi(t):t\in[0,1]\}$ and
the smallest closed interval containing 
$\{p_k\zeta(t):t\in[0,1]\}$.)
In particular,
\[
H\le\sum_{k=1}^d\Big\{\Big|\max_{[0,1]}p_k\xi-\max_{[0,1]}p_k\zeta\Big|\vee
\Big|\min_{[0,1]}p_k\xi-\min_{[0,1]}p_k\zeta\Big|\Big\}
\le d_\tD(\xi,\zeta).
\]
Again,
\[
d_{E,[0,t_{n,N_n}]}(\pi_E \tU_n,\pi_E \chi U_n)
= b_n^{-1}\max_{0\le j\le n}H\circ f^j
\le  b_n^{-1}\max_{0\le j\le n}d_\tD(\xi,\zeta)\circ f^j,
\]
completing the proof of the claim in the case $\disc_{U_n}=\tcS_n$ for all $n$. 

In general, there is the possibility that $t_{n,j+1}\in\tcS_n\setminus \disc_{U_n}$.
On the interval $(t_{n,j},t_{n,j+1}]$, it remains the case that
$\pi_\tD \tU_n$ is the concatenation of the constant function $U_n(t_{n,j})$ with $U_n(t_{n,j})+b_n^{-1} \xi(t) \circ f^j$, while
$\pi_\tD \chi U_n$ is simply the 
constant function $U_n(t_{n,j})$. The latter is equivalent to
the concatenation of $U_n(t_{n,j})$ with 
$U_n(t_{n,j})$. 
But the fact that $t_{n,j+1}$ is not a discontinuity point of $U_n$ means that $V=0$ and hence (by definition) $\zeta(t)\circ f^j=0$.
Hence
$\pi_\tD \chi U_n$ is still the
concatenation of 
$U_n(t_{n,j})$ with
$U_n(t_{n,j})+b_n^{-1}\zeta(t)\circ f^j$. 
\end{proof}

\end{document}
